\scanpage{128}
\mspage{63}
Il faut introduire encore les espaces analytiques
de drapeaux $D_{d_0,d_1\mathord{-\, -},d_r}$ pour $d_0\geq \unsure{d_2|d_1}\geq\mathord{\cdot\,-}\geq d_r\geq0$
leur somme est notée $D_{\Delta_r}$, où $\Delta_r$ désigne
maintenant l'ens. ordonné
\[
\tag{203}\Delta_r=[0,r]\cap\mathbb N
\]
ou simplexe ordonné type de dim $r$. Cette notation
$D_{\Delta_r}$ est compatible avec la notation $D_J$ pour $J$
ens. fini ordonné ayant un plus petit élément.
Alors $D_{I_r}$ dépend fonctoriellement de façon
contravariante en $I_r$, pour les applications
croissantes, entre simplexes, donc on a un
espace analytique $1/2$-simplicial, qu'on désigne par
$\mathcal D_*$, donc
\[
\tag{204}
\underset{\inserted{\substack{\text{ou}\\D_{\Delta_r}}}}{\mathcal D_r}
=\overset{\inserted{D_{\Delta_r}}}{\deleted{\unsure{D_{\Delta_I}|D_{\Delta_r}|D_{I_r}}}}
=\coprod_{\substack{(d_0,\mathord{-\,\cdot},d_r)\in\mathbb N^r\\d_0>d_1\mathord{\cdot\,\cdot\,-}>d_r}}
D_{d_0,d_1,\mathord{-},d_r}
\]
\deleted{$\unsure{D_{d_0,d_r}|D_{\Delta_r,d_r}}$}
La composante $D_{d_*}=D_{d_0,d_1,\unsure{,}\mathord{-},d_r}$ est
un \unsure{revêt} étale fini de $D_{d_0}$, s'identifiant
\unsure{:|à}
\[
\tag{205}
D_{d_*}=D^{!}_{d_0}\wedge_{(\mathfrak S_{d_0})_{D_{d_0}}}
\underbrace{\mathcal P_{d_*}\bigl(\deleted{\unsure{[1,d_0]}}I_{d_0}\bigr)}_{
\substack{\text{ens des ``\deleted{\illegible{1}}\inserted{drapeaux} de}\\
\text{type }d_*\text{'' de l'ens fini}\\
I_{d_0}=[1,d_0]\cap\mathbb N\quad\text{\tikz[x=.07cm,y=.07cm,baseline=0pt]{\draw[line width=.3pt] (0,0) .. controls (2.1,.4) and (2.9,.4) .. (3,2);}}}}
\]
on désigne par $D^*_{d_0,d_1,\mathord{-}d_r}$ l'image inverse
de $D^*_{d_0}$, qui est donc le complémentaire
de l'image inverse $\Delta_{d_0-d_r}$ du diviseur
$\Theta_{d_0}$ sur $D_{d_0}$. On peut bien sûr plus simple-
\margin{127}
